\begin{statementbox}
\textbf{\textcolor{CStatement}{这条结论解决什么问题？}}

遇到 $y=f(g(x))$ 这样“函数里面还有函数”的式子，怎样不用把式子展开，
也不必一上来求导，就能判断它随 $x$ 的增大是增大还是减小？
答案是：\textbf{先看里面怎样变化，再看外面会不会把大小关系反过来。}

\medskip
\textbf{\textcolor{CStatement}{从哪里来？}}

例如先算 $10-x$，再把结果代入 $10-t$：
\[
 x\longrightarrow 10-x\longrightarrow 10-(10-x)=x.
\]
第一次让数“越往右越小”，第二次又把大小关系反过来，
最后反而随 $x$ 增大而增大。\textbf{两次反向，方向恢复；一次反向，方向相反。}
这就是“同增异减”产生的原因，而不只是一个要背的口诀。

\medskip

\begin{center}\begin{tikzpicture}[font=\small,>=stealth,line join=round]
\draw[gray!40,rounded corners=3pt] (0,0) rectangle (14.2,9.0);
\node at (7.1,8.4) {先减一次，再减一次：大小顺序恢复};
\draw[->,gray] (1.0,7.2)--(7.5,7.2) node[right]{$x$};
\fill[blue] (2.3200000000000003,7.2) circle(2.2pt);
\node[above left=7pt,fill=white,inner sep=1pt] at (2.3200000000000003,7.2) {$2$};
\draw[blue,dashed,thick,->,shorten <=3pt,shorten >=3pt] (2.3200000000000003,7.2)--(5.98,4.7);
\fill[orange!85!black] (4.15,7.2) circle(2.2pt);
\node[above right=7pt,fill=white,inner sep=1pt] at (4.15,7.2) {$5$};
\draw[orange!85!black,dashed,thick,->,shorten <=3pt,shorten >=3pt] (4.15,7.2)--(4.15,4.7);
\draw[->,gray] (1.0,4.7)--(7.5,4.7) node[right]{$t$};
\fill[blue] (5.98,4.7) circle(2.2pt);
\node[above right=7pt,fill=white,inner sep=1pt] at (5.98,4.7) {$8$};
\draw[blue,dashed,thick,->,shorten <=3pt,shorten >=3pt] (5.98,4.7)--(2.3200000000000003,2.2);
\fill[orange!85!black] (4.15,4.7) circle(2.2pt);
\node[above left=7pt,fill=white,inner sep=1pt] at (4.15,4.7) {$5$};
\draw[orange!85!black,dashed,thick,->,shorten <=3pt,shorten >=3pt] (4.15,4.7)--(4.15,2.2);
\draw[->,gray] (1.0,2.2)--(7.5,2.2) node[right]{$y$};
\fill[blue] (2.3200000000000003,2.2) circle(2.2pt);
\node[below left=7pt,fill=white,inner sep=1pt] at (2.3200000000000003,2.2) {$2$};
\fill[orange!85!black] (4.15,2.2) circle(2.2pt);
\node[below right=7pt,fill=white,inner sep=1pt] at (4.15,2.2) {$5$};
\node[align=left,anchor=west] at (8.4,4.7){$g(x)=10-x$：减\\[12pt]$f(t)=10-t$：减\\[14pt]$H(x)=x$：增\\[14pt]$2<5$\\[9pt]$8>5$\\[9pt]$2<5$};
\node[text width=12.8cm,align=center] at (7.1,.65) {蓝点与橙点各自代表同一个输入；各行数轴向右增大，连接线追踪实际数值。};
\end{tikzpicture}\end{center}

\end{statementbox}
\clearpage
\begin{statementbox}
\textbf{\textcolor{CStatement}{适用条件（先检查这三项）}}
\begin{enumerate}[leftmargin=*, itemsep=0.25em]
\item 先选定一个自变量区间 $I$，使 $g(x)$ 在这段上始终递增，或始终递减。
\item 对 $x\in I$，$g(x)$ 得到的每个值，都必须能代入外层 $f$。
\item \textbf{最容易漏：}外层 $f(t)$ 对这些实际得到的 $t=g(x)$，也必须始终递增，或始终递减；仅仅“有定义”并不够。
\end{enumerate}

\textbf{\textcolor{CStatement}{核心结论：同增异减}}

在以上条件下，若两层均为\textbf{严格}递增或严格递减，则
\[
\begin{array}{c|c|c}
\text{内层 }g & \text{外层 }f & \text{整体 }f(g(x))\\ \hline
\text{增} & \text{增} & \text{增}\\
\text{减} & \text{减} & \text{增}\\
\text{增} & \text{减} & \text{减}\\
\text{减} & \text{增} & \text{减}
\end{array}
\]
其中“外层增或减”必须针对\textbf{内层实际输出的数值范围}来判断，
不是只看 $f$ 在某个无关区间上的增减。

\medskip
\textbf{\textcolor{CStatement}{多层推广：数一数反向的次数}}

三层、四层也一样：每经过一层\textbf{递减}，大小关系就反转一次。
只要每层在上一层实际给它的数值范围内都保持同一种增减方向，
\textbf{递减层数是偶数，整体递增；是奇数，整体递减。}
也可把增记为 $+1$、减记为 $-1$，逐层相乘：
\[
 s_{\text{整体}}=s_1s_2\cdots s_n,\qquad s_i\in\{+1,-1\}.
\]

\textbf{边界提醒：}如果允许“保持不变”这样的情况，只能相应判断\emph{不减或不增}，
不能保证\emph{严格}增减；若有某层在有关范围内是常数，后续各层都有定义时，整体就是常数。
\end{statementbox}
